\section{Performance and Accuracy}

\subsection{Measurement boundary}

All headline backward rows are causal B1 with $D_{QK}=192$ and $D_V=128$ on
GPU 0.  The copied BF16 V382 mode and low-precision mode run through the same
extension and scheduler family.  Table~\ref{tab:direct-bf16-results} uses five
warmups and 21 rotated CUDA-event samples with seed 2026081201.  Operand
preparation is outside the standalone backward timing; projection tables state
explicitly whether BF16-to-NVFP4 activation preparation is charged.

These are implementation medians, not peak-throughput estimates.  The result
does not include an optimizer step, projection weight-gradient GEMM, or full
training iteration.

\subsection{Standalone backward}

Reducing dQ in BF16 and returning it directly gives the strongest deployable
backward endpoint:

\begin{table}[htbp]
\centering
\caption{Retained adaptive backward versus the copied BF16 control.  Lower is
better.}
\label{tab:direct-bf16-results}
\small
\begin{tabular}{rrr r r}
\toprule
S & H & BF16 control (ms) & adaptive BF16 dQ (ms) & improvement \\
\midrule
4096  & 24 & 0.704480 & 0.624832 & 11.31\% \\
8192  & 8  & 0.868864 & 0.746432 & 14.09\% \\
8192  & 24 & 2.207200 & 1.888768 & 14.43\% \\
8192  & 64 & 5.560512 & 4.761888 & 14.36\% \\
16384 & 24 & 7.743616 & 6.471232 & 16.43\% \\
16384 & 64 & 20.273727 & 17.068480 & 15.81\% \\
\bottomrule
\end{tabular}
\end{table}

The FP32-return API uses the same BF16 reduction and adds one widening pass.
Skipping that pass saves 2.0--5.0\% relative to the FP32-return adaptive route.
Changing the reduction payload itself from FP32 to BF16 accounts for another
0.5--5.2\%, with the benefit increasing at long sequence.

\subsection{Calibrated and outlier accuracy}

Table~\ref{tab:direct-bf16-quality} compares direct BF16 outputs to the copied
BF16 control.  Values are component cosine, not an aggregate dominated by dV.

\begin{table}[htbp]
\centering
\caption{Calibrated component cosine for the direct-BF16 endpoint.}
\label{tab:direct-bf16-quality}
\small
\begin{tabular}{rr rrr}
\toprule
S & H & dQ & dK & dV \\
\midrule
4096  & 24 & 0.991447 & 0.991379 & 0.999356 \\
8192  & 8  & 0.991533 & 0.991557 & 0.999378 \\
8192  & 24 & 0.991435 & 0.991484 & 0.999360 \\
8192  & 64 & 0.991478 & 0.991487 & 0.999354 \\
16384 & 24 & 0.991560 & 0.991526 & 0.999355 \\
16384 & 64 & 0.991447 & 0.991452 & 0.999345 \\
\bottomrule
\end{tabular}
\end{table}

The robust scale rule in Equation~\ref{eq:robust-clip} matters outside the
calibration distribution.  On a one-percent sparse Q/K outlier test:

\begin{table}[htbp]
\centering
\caption{Sparse-outlier accuracy of the retained adaptive route.}
\label{tab:outlier-quality}
\small
\begin{tabular}{rr rrr}
\toprule
S & H & dQ cosine & dK cosine & dV cosine \\
\midrule
4096 & 24 & 0.962695 & 0.962139 & 0.999070 \\
8192 & 8  & 0.961676 & 0.963121 & 0.998915 \\
\bottomrule
\end{tabular}
\end{table}

The earlier global-amax adaptive rule reached only about 0.753/0.757 dQ/dK
cosine at S4096/H24.  The robust per-head rule improves the outlier regime
without changing the three score instructions, TMEM map, or calibrated result.
It still clips information and is not a proof of training convergence.

\subsection{Learned projection publication}

Table~\ref{tab:projection-charged} charges the current standalone conversion
of BF16 X into the NVFP4 activation operand.  NVFP4 weights are treated as
persistent and prepared once.  The low-precision column uses the automatic
publication policy in Equation~\ref{eq:projection-dispatch}.

\begin{table}[htbp]
\centering
\caption{Q/K learned-projection plus backward-layout production.  The
BF16-input column includes current activation preparation.}
\label{tab:projection-charged}
\small
\begin{tabular}{l rrr}
\toprule
Shape & BF16 GEMM + pack (ms) & BF16 input to NVFP4 auto (ms) & speedup \\
\midrule
S8192/H8/K1024  & 0.163712 & 0.142768 & 1.147$\times$ \\
S4096/H24/K3072 & 0.326704 & 0.206096 & 1.585$\times$ \\
S4096/H64/K8192 & 1.553184 & 0.515616 & 3.012$\times$ \\
\bottomrule
\end{tabular}
\end{table}

If X already arrives in NVFP4, the corresponding hot-stage speedups are
1.473$\times$, 1.933$\times$, and 3.342$\times$.  Q/K cosine against the
BF16 learned projection is approximately 0.99097--0.99098.  Fused versus
separate E2M1 publication is bit-exact.

The dispatcher removes the shallow-K epilogue regression:

\begin{table}[htbp]
\centering
\caption{Automatic projection publication versus forcing the other route.}
\label{tab:projection-dispatch-results}
\small
\begin{tabular}{l l rrr}
\toprule
Shape & auto route & auto (ms) & forced other (ms) & advantage \\
\midrule
S8192/H8/K1024  & separate & 0.111024 & 0.122368 & 1.102$\times$ \\
S4096/H24/K3072 & fused    & 0.169984 & 0.181712 & 1.069$\times$ \\
S4096/H64/K8192 & fused    & 0.462560 & 0.626592 & 1.355$\times$ \\
\bottomrule
\end{tabular}
\end{table}

The dominant gain is real NVFP4 tensor work and reduced A/B traffic.  At deep
K, the fused epilogue adds a movement win by avoiding the BF16 Q/K reread and
standalone transpose.  At K1024, the wider parallel packer is faster than
placing all publication on the persistent GEMM's one consumer warpgroup.
These two tables describe the legacy Q/K-only API.

The unified QKV producer adds projection-native MXFP4 V and allows all three
BF16 stores to be suppressed.  Table~\ref{tab:unified-qkv-projection} compares
it with three BF16 GEMMs followed by the standalone Q/K packer.  Input and
weight NVFP4 operands are already resident in all columns so that the table
isolates projection and publication policy.

\begin{table}[htbp]
\centering
\caption{Unified QKV projection and native forward/backward publication.}
\label{tab:unified-qkv-projection}
\small
\begin{tabular}{l rrr rr}
\toprule
Shape & BF16 + QK pack & unified + BF16 & speedup & unified no BF16 & speedup \\
 & (ms) & (ms) & & (ms) & \\
\midrule
S2048/H8/K1024   & 0.149264 & 0.088688 & 1.683$\times$ & 0.077440 & 1.927$\times$ \\
S4096/H24/K4096  & 0.463744 & 0.328816 & 1.410$\times$ & 0.289552 & 1.602$\times$ \\
\bottomrule
\end{tabular}
\end{table}

Projection Q/K/V cosine is 0.99096--0.99098.  Consuming the native operands
in causal forward gives output cosine 0.990722 at S2048/H8 and 0.991317 at
S4096/H24.  All four backward Q/K layouts are bit-exact against their
standalone publishers.  Under the retained 2-D policy, the forward- and
backward-oriented V payloads are exact transposes and their repeated scale
pages select the same $32\times32$ tile scale.  The legacy standalone V
publishers select independent $1\times32$ scales, so they are intentionally
no longer a byte-level reference.  The forward Q/K tensors and backward
compact tensors alias the same allocation.

The storage reduction is material at long sequence.  Sharing the compact Q/K
payload avoids $BHS\cdot192$ bytes: 12 MiB at S8192/H8 and 18 MiB at
S4096/H24.  Suppressing BF16 Q/K/V publication avoids a further 64 MiB and
96 MiB respectively.  The projection producer also reduces its static shared
scratch by 2,112 bytes, from 10,704 to 8,592 bytes.  This scratch saving does
not change the already near-full shared-memory allocation of the backward
kernel.

\subsection{Two-dimensional scale-contract audit}

The composed D64 path was re-audited with the representation rules in
Equations~\ref{eq:nvfp4-projection-weight-2d} and
\ref{eq:projection-native-v}.  Packing a synthetic weight and independently
packing its transpose produces exactly the transposed E2M1 codes, matching
E4M3 local scales, and the same global scale under the $16\times16$ policy.
The $1\times16$ policy does not provide that invariant.

Sharing scales over two dimensions is stricter and therefore slightly less
accurate for a single isolated tensor.  On the fixed D64 projection audit,
learned-projection cosine changes from 0.990916 with $1\times16$ weights to
0.989259 with $16\times16$ weights; relative L2 changes from 0.13483 to
0.14630.  Decoded V cosine changes from 0.993560 with independently selected
$1\times32$ row/column scales to 0.993031 with one $32\times32$ scale;
relative L2 changes from 0.11332 to 0.11789.  The retained reason for the
two-dimensional policy is not a better one-shot cosine: it is transpose
consistency when fprop and bprop actually consume the two orientations of the
same low-precision representation.  That benefit is only partially active in
the selected hybrid.  Output projection does consume both orientations, but
QKV projection dgrad remains BF16 and attention backward consumes FP8 V.
Consequently, applying one 2-D switch to every learned weight and to V also
charges coarser forward quantization where no transposed FP4 consumer reuses
it.  The policy remains useful for the all-FP4 ladder; the hybrid requires
separate QKV/output-weight switches rather than a blanket 2-D default.

The E8M0 exponent selector must likewise follow the scale geometry.  On the
fixed causal-V activation audit, the $1\times32$ policy minimizes attention
relative L2 at 0.11409 when it rounds up at normalized amax 1.20; ordinary
nearest-power rounding (the 1.50 midpoint) gives 0.12818.  For a shared
$32\times32$ tile the ordering reverses: nearest-power rounding gives 0.12439,
versus 0.12745 with the 1.20 rule.  The retained publisher therefore uses a
BF16-exact 1.203125 cutoff for $1\times32$ and ties-to-even nearest-power
rounding for $32\times32$.  A staged binary matched the independent payload
and scale emulator bit-for-bit for both geometries, with and without causal
interleaving, and passed the unified projection-layout validator.  These are
one-activation representation checks, not convergence measurements.  In the
matched sixteen-layer diagnostic, the new selector improves first-layer
attention-branch relative L2 from 0.37686 to 0.37438, but the final sampled
logit changes are mixed and negligible: cosine changes from 0.68707 to
0.68597 while relative L2 changes from 0.79662 to 0.79494.  Selector rounding
therefore does not explain the much larger optimization-proxy regression.

The implementation cost is below whole-step measurement noise.  In matched
16-layer S4096/H32 GQA runs, the low-precision step measured 76.929~ms with
legacy one-dimensional scales and 76.931~ms with two-dimensional scales.  The
two-dimensional run used independent Q/K FP4 multipliers 2.25 and 2.0 and
unit field gains.  Against an 83.304-ms BF16 step in that run, it measured
1.083$\times$ end-to-end speedup.  This short synthetic-token comparison is a
systems measurement; its 16-layer sampled-logit cosine of 0.6989 is not a
convergence result.

The backward scale audit also found and removed a stale common gain.  The
previous MXFP4-PV comparison silently inferred 0.632 from its forward softmax
mode, yielding initial Q/K weight-gradient norm ratios near 0.593 and a V
ratio near 0.425.  With Equation~\ref{eq:d64-field-local-gradient-decode} and
unit gains, the corresponding Q/K/V ratios are approximately
0.954/0.937/0.961.  Directional error remains after sixteen approximate
layers, however, and norm parity does not imply a more accurate gradient.  If
$r=\lVert\widehat g\rVert/\lVert g\rVert$ and $c$ is the cosine to the BF16
gradient, applying gain $a$ gives
\begin{equation}
  \frac{\lVert a\widehat g-g\rVert_2}{\lVert g\rVert_2}
  = \sqrt{1+(ar)^2-2arc}.
  \label{eq:gain-direction-error}
\end{equation}
For $c\approx0.3$, forcing $ar\approx1$ raises relative error to about 1.18.
Thus the unit-gain endpoint removed a norm confound but amplified a poorly
aligned direction.  It is an audit point, not a calibrated training policy.
The optional NVFP4 QKV-dgrad consumer now folds the same three decode factors
into its E4M3 metadata.  It restores the expected norms but is not selected:
its measured backward is 46.11~ms versus 43.89~ms for the BF16 projection
dgrad route.

No random Hadamard transform or stochastic rounding is charged in the
selected route.  Forward activations use $1\times16$ RNE, learned weights use
$16\times16$ RNE, and QKV projection dgrad plus both weight-gradient GEMMs
remain BF16.  RHT/SR becomes relevant only if those gradient operands are
actually retained in FP4; adding it around the slower diagnostic NVFP4 dgrad
would be fake quantization overhead rather than a numerical fix for the
winner.

A controlled 24-round comparison then put both PV routes on this corrected
contract.  Table~\ref{tab:d64-2d-route-comparison} reports medians for the
synthetic 16-layer S4096 Llama-1.2B training proxy.

\begin{table}[htbp]
\centering
\caption{End-to-end D64 route comparison with 2-D learned weights, 2-D MXFP4
V, independent Q/K multipliers, and unit Q/K/V backward gains.}
\label{tab:d64-2d-route-comparison}
\small
\begin{tabular}{l rrr r r}
\toprule
Route & forward & backward & optimizer & step & speedup \\
 & (ms) & (ms) & (ms) & (ms) & vs BF16 \\
\midrule
BF16 CuTE                 & 27.718 & 45.596 & 10.726 & 84.050 & 1.000$\times$ \\
NVFP4 QK + MXFP4 PV       & 22.466 & 44.239 & 10.732 & 77.428 & 1.086$\times$ \\
NVFP4 QK + exact FP8 PV   & 24.884 & 44.233 & 10.733 & 79.860 & 1.052$\times$ \\
\bottomrule
\end{tabular}
\end{table}

The two low-precision backward times are indistinguishable at this scale; the
MXFP4 route's additional 2.43-ms whole-step gain comes from forward PV.  The
corresponding sampled-logit cosines are 0.69348 for MXFP4 PV and 0.77821 for
exact FP8 PV.  Q/K/V first-layer weight-gradient cosine is
0.27691/0.32003/0.39347 for MXFP4 and
0.44378/0.45383/0.68284 for FP8, while Q/K gradient norm ratios remain near
0.94--0.96 in both routes.  Therefore the corrected difference is primarily
directional PV error, not the removed common-gain bug.  The routes share the
same QK policy only at the first layer: their distinct PV outputs perturb the
residual stream, so subsequent layers quantize different hidden states.

All routes remained finite in the short 24-step optimization proxy.  The
last-cycle median losses were 4.281 for BF16, 5.797 for MXFP4 PV, and 5.344
for exact FP8 PV.  A later audit found that the compile warmup performed one
unreported optimizer update, so these values must not be used to select a
training policy.  The speed measurements remain systems measurements, but
the apparent loss ordering motivated the controlled ablation below.

\subsection{Post-hoc loss-regression ablation}
\label{sec:loss-regression-ablation}

The endpoint above changed four policies together: learned weights from 1-D
to 2-D scales, V from 1-D to 2-D scales, Q/K multipliers from 1.75/1.75 to
2.25/2.0, and the inferred backward gain from 0.632 to 1.0.  An initial
eight-run factorial held seeds, token batches, route order, and optimizer
settings fixed.  Before discovery of the hidden warmup update, that factorial
attributed 0.596 of the 0.691 mean MX--BF16 loss-gap regression (86.3\%) and
0.941 of the 1.037 last-four regression (90.8\%) to the gain change.  Those
endpoint losses are not themselves valid policy evidence, but the factor
ranking motivated the corrected controls below.  Changing Q/K multipliers
alone is neutral within run noise.  Two-dimensional learned weights reduce
initial sampled-logit cosine by about 0.029 in both PV routes and modestly
worsen the loss proxy; two-dimensional V has a much smaller deterministic
effect.  Gain changes initial logits by exactly zero.

The mechanism follows Equation~\ref{eq:gain-direction-error}.  In the
sixteen-layer endpoint, Q/K gradient cosine is only about 0.28--0.32 for MX
and 0.44--0.45 for FP8.  Raising MX Q/K/V norm ratios from roughly
0.59/0.60/0.44 to 0.94/0.95/0.75 therefore increases relative L2 even though
the norms look healthier.  By contrast, the isolated one-layer audit gives
Q/K/V weight-gradient cosine 0.89/0.89/0.91 for MX and
0.93/0.92/0.93 for FP8.  The local stitch operation is bit-exact for Q, K,
and V, while the attention kernel itself gives dQ/dK cosine near 0.997 and dV
near 0.9995.  The poor sixteen-layer direction is therefore accumulated
forward-trajectory error, not a local backward layout failure.

The warmup was changed to allocate and compile fused AdamW state at zero
learning rate, reset its moments and step counters, and leave parameters
unchanged.  Under this corrected harness, audited initial loss equals recorded
round-zero loss exactly.  Table~\ref{tab:clean-gain-sweep} shows that the
four-batch memorization proxy improves monotonically as the common gain tends
toward zero.

\begin{table}[htbp]
\centering
\caption{Corrected 24-step repeated-batch proxy.  Entries are mean relative
loss differences to each run's BF16 route; negative is lower on this proxy.}
\label{tab:clean-gain-sweep}
\small
\begin{tabular}{r rr rr}
\toprule
Common gain & MX mean & MX last four & FP8 mean & FP8 last four \\
\midrule
0.02 & -25.29\% & -45.93\% & -37.29\% & -74.12\% \\
0.05 & -23.05\% & -42.56\% & -34.27\% & -67.30\% \\
0.10 & -19.40\% & -35.77\% & -28.90\% & -58.42\% \\
0.20 & -14.50\% & -25.35\% & -23.36\% & -46.68\% \\
0.40 &  -7.36\% & -11.23\% & -15.22\% & -33.56\% \\
0.50 &  -4.35\% &  -2.06\% & -11.85\% & -26.11\% \\
0.632 &  2.05\% & 10.91\% & -5.16\% & -9.49\% \\
1.00 & 14.27\% & 39.15\% &  9.58\% & 26.37\% \\
\bottomrule
\end{tabular}
\end{table}

The corrected matched endpoints independently reproduce the original factor
ranking.  Moving from 0.632 to 1.0 worsens the mean and last-four relative
gaps by 12.22 and 28.25 percentage points for MXFP4 PV, and by 14.74 and
35.86 points for FP8 PV.  The BF16 last-cycle median differs by only 0.0078
between those separately scheduled runs.

This monotone result is a harness pathology, not a low-precision win.  A gain
of 0.02 nearly freezes Q/K/V projection-weight gradients and the attention
contribution to upstream $dX$, while residual, MLP, embedding, and output
projection paths continue to memorize four repeated random sequences.  The
proxy therefore rewards routing learning around attention and cannot select
a production gain.

A unique-batch control confirms that interpretation.  When the same 24 steps
consume 24 independently generated random sequences, every route begins and
ends at BF16-rounded loss 12.1875 and mean loss remains approximately 12.18.
Changing the gain from 0.02 to 1.0 then changes the mean relative route gap by
only 0.0001 percentage points for MXFP4 PV and 0.0427 points for FP8 PV.  On
the four-batch cycling proxy, the same endpoint comparison changes the mean
gap by 39.56 and 46.87 points.  The dramatic gain curve is therefore a
repeated-example memorization effect, not evidence of model convergence or
generalization.

A field ablation makes the failure more specific.  Starting from all fields
at 0.02, raising Q or K to 0.5 slightly improves both routes, and raising both
improves the mean absolute loss gap by 0.163 (MX) and 0.188 (FP8).  Raising V
alone worsens it by 1.513 and 1.850 respectively.  A complete $2\times2$
V-gain factorial separates V's projection-weight update from its contribution
to upstream $dX$.  Relative to the low/low corner, raising only V-to-$dX$
worsens the mean absolute loss gap by 1.513 (MX) and 1.888 (FP8), whereas
raising only the V-weight gain changes it by just 0.051 and $-0.020$.
The high/high corner reproduces the V-field regression.  This is a
route-common property of this proxy---the FP8 effect is larger---rather than
evidence for an MXFP4-specific V-layout defect.  The post-attention stitch was
independently checked bit-exactly, and isolated dV cosine exceeds 0.9993, which
rules out a gross transposition or decode-scale bug.  The factorial establishes
that V-to-$dX$ controls this proxy response; it does not establish that this
path uniquely controls real-training convergence.

The actionable conclusion is twofold.  First, no gain from this synthetic
curve should be promoted; gain and learning rate must be co-tuned on real,
non-repeated tokens with held-out loss.  Second, numerical work should target
forward trajectory fidelity---especially the probability/V path---rather than
normalizing already misaligned full-model gradients.

\subsection{RoPE-aware projection and composed decoder boundary}

Table~\ref{tab:rope-projection-components} isolates full-head pair-native
RoPE at S4096/H24/K3072.  ``Separate'' runs the same projection and then a
strong in-place BF16 Q/K rotary kernel.  The retained fused implementation
uses packed BF16 arithmetic in the projection epilogue.

\begin{table}[htbp]
\centering
\caption{Median component timing for pair-native RoPE publication.}
\label{tab:rope-projection-components}
\small
\begin{tabular}{l r}
\toprule
Operation & median (ms) \\
\midrule
unrotated unified projection & 0.341280 \\
projection with fused packed-BF16 RoPE & 0.432384 \\
projection then standalone RoPE & 0.448512 \\
standalone forward RoPE only & 0.114304 \\
standalone inverse-RoPE gradient only & 0.114272 \\
\bottomrule
\end{tabular}
\end{table}

Fusion saves 16.128~$\mu$s against the separate forward transform, but it
does not make RoPE free: the rotation extends the projection consumer path by
91.104~$\mu$s.  The packed-BF16 result has Q/K cosine 0.9999978 and relative
L2 approximately 0.00210 against an FP32-FMA epilogue reference; its maximum
absolute difference is 0.00390625.  All four E2M1 layouts are byte-exact
against standalone publication of the same rotated BF16 tensors, and V is
bit-exact.  End-to-end Q/K projection cosine against dense standard-order
RoPE is approximately 0.990966, so the incremental rotary rounding is small
relative to the native NVFP4 projection error.

The inverse-gradient handoff has a different ownership problem.  The generic
NVFP4 quantizer assigns one complete 128-column row to each thread, whereas
the standalone inverse transform assigns adjacent rotary pairs across a warp.
Table~\ref{tab:inverse-rope-nvfp4-handoff} measures the retained direct-table
specialization at the same S4096/H24 shape.  Payloads, E4M3 scale pages, and
the optional BF16 republication are bit-exact with the separate reference.

\begin{table}[htbp]
\centering
\caption{Median inverse-RoPE and delayed-scale NVFP4 handoff timing at
S4096/H24.}
\label{tab:inverse-rope-nvfp4-handoff}
\small
\begin{tabular}{l r}
\toprule
Operation & median (ms) \\
\midrule
standalone inverse RoPE & 0.114592 \\
NVFP4 pack with precomputed global scale & 0.059968 \\
separate inverse RoPE then NVFP4 pack & 0.159712 \\
fused handoff, no BF16 republication & 0.156992 \\
fused handoff with BF16 republication & 0.156320 \\
\bottomrule
\end{tabular}
\end{table}

The no-publication path saves 2.720~$\mu$s, or 1.70\%, and compiles at 122
registers with no spills.  The publication form compiles at 153 registers
without spills and saves 3.392~$\mu$s, or 2.12\%; its asynchronous BF16 store
is overlapped by the packing tail.  Cooperative table staging was also
tested: despite coalescing the cos/sin reads, its longer live range raised the
kernels to 182--250 registers and made them substantially slower.  Inverse
RoPE is therefore not the missing source of a forward-sized end-to-end gain;
eliminating the BF16 gradient boundary itself remains necessary.

A rotated 31-sample comparison in the composed harness confirms the system
effect rather than extrapolating it from component times:

\begin{table}[htbp]
\centering
\caption{Fused inverse-RoPE handoff at S4096/H24/K4096 after seven warmups.}
\label{tab:inverse-rope-e2e-ab}
\small
\begin{tabular}{l rrrr}
\toprule
Scope & BF16 (ms) & separate (ms) & fused (ms) & fused vs BF16 \\
\midrule
activation gradients & 2.188992 & 1.701568 & 1.693664 & 1.292$\times$ \\
full training & 2.620640 & 2.138720 & 2.125952 & 1.233$\times$ \\
\bottomrule
\end{tabular}
\end{table}

Relative to the same low-precision path with a separate inverse transform,
fusion reduces activation-scope time by 7.904~$\mu$s (0.46\%) and full
training time by 12.768~$\mu$s (0.60\%).  The fused BF16 republication makes
the projection weight-gradient result bit-identical to the separate path, so
the training row does not omit a required consumer.

Table~\ref{tab:rope-e2e} measures the dense causal decoder-attention boundary
with QKV projection, fused RoPE, retained FP4 forward, output projection,
output-projection dgrad, attention backward, inverse RoPE, and QKV projection
dgrad.  ``Activation'' omits projection weight gradients; ``training'' adds
both projection weight-gradient GEMMs.  Persistent weight operands, rotary
tables, pair-native weight order, adaptive head records, and delayed global
scales are cached outside timing.  The selected column is the upstream-native
delayed-scale NVFP4 route before the inverse-handoff fusion isolated in
Table~\ref{tab:inverse-rope-e2e-ab}.

\begin{table}[htbp]
\centering
\caption{RoPE-aware composed boundary on GPU 0.  Times are medians over 13
samples after four warmups.}
\label{tab:rope-e2e}
\scriptsize
\begin{tabular}{rrr l rrr}
\toprule
S & H & K & scope & BF16 (ms) & low precision (ms) & speedup \\
\midrule
4096  & 24 & 4096 & activation & 2.195808 & 1.701600 & 1.290$\times$ \\
4096  & 24 & 4096 & training   & 2.627328 & 2.144448 & 1.225$\times$ \\
8192  & 24 & 4096 & activation & 5.037664 & 4.206208 & 1.198$\times$ \\
8192  & 24 & 4096 & training   & 5.875584 & 5.039552 & 1.166$\times$ \\
4096  & 64 & 8192 & activation & 7.119584 & 4.401792 & 1.617$\times$ \\
4096  & 64 & 8192 & training   & 9.208512 & 6.464608 & 1.424$\times$ \\
8192  & 64 & 8192 & activation & 16.682880 & 11.233536 & 1.485$\times$ \\
8192  & 64 & 8192 & training   & 20.868896 & 15.431264 & 1.352$\times$ \\
16384 & 64 & 8192 & activation & 44.307362 & 33.336990 & 1.329$\times$ \\
16384 & 64 & 8192 & training   & 52.857922 & 42.032127 & 1.258$\times$ \\
\bottomrule
\end{tabular}
\end{table}

Across these rows, output-projection $Y$ cosine is 0.97241--0.97397, dX is
0.98115--0.98146, dQ is 0.95368--0.96084, dK is 0.95348--0.95956, and dV is
0.99006--0.99045.  The boundary is intentionally not advertised as a stock
Llama drop-in: it retains $D_{QK}=192$, equal Q/K/V head counts, and lacks
grouped-query attention and native $D_{QK}=128$ kernels.  It is nevertheless
a RoPE-correct dense decoder attention measurement rather than the earlier
RoPE-free proxy.  The separate inverse transform costs about 0.115~ms at the
component shape.  The fused handoff recovers only 2.720--3.392~$\mu$s within
the isolated producer because its row-owned quantizer mapping makes
rotary-table access strided.  The composed gain is 0.46--0.60\% at
S4096/H24/K4096.  A larger gain requires the projection consumer to accept
the completed reduction tile directly, not merely to fuse arithmetic into a
later global-memory quantizer.

\subsection{Stacked QKV-gradient handoff}

Table~\ref{tab:stacked-qkv-handoff} compares the two variants in
Equations~\ref{eq:stacked-qkv-gradient-row}--\ref{eq:stacked-qkv-gradient-projection}
against the previously selected interleaved-gradient route.  All three rows
use the retained FP4 forward and FP4+FP8 backward.  Times are complete
activation-gradient boundaries, not isolated packers.

\begin{table}[htbp]
\centering
\caption{Stacked QKV-gradient publication at S4096/H24/K3072.  Medians use
21 rotated samples after eight warmups.}
\label{tab:stacked-qkv-handoff}
\small
\begin{tabular}{l rr}
\toprule
Route & time (ms) & result versus selected \\
\midrule
interleaved QKV gradient + fused inverse-RoPE pack & 1.578336 & baseline \\
unsignaled one-lane stacked QKV pack + projection & 1.566048 & 0.78\% faster \\
signaled two-lane hierarchical pack + projection & 1.635936 & 3.65\% slower \\
\bottomrule
\end{tabular}
\end{table}

The one-lane output has dX cosine 0.981765 against the complete BF16 chain;
against the selected low-precision dX its cosine is 0.99999994 and relative
L2 is $6.45\times10^{-5}$.  dK and dV are bit-identical to the selected
attention outputs before the common projection quantization.  The two-lane
form is numerically valid but loses because tile-ready signaling adds
approximately 37~$\mu$s even in a one-lane control and roughly 48~$\mu$s for
the actual hierarchy.  Its arithmetic savings do not repay that perturbation.
The unsignaled form is therefore retained only for activation dgrad.  Full
training retains the interleaved BF16 publication because the QKV
weight-gradient GEMM requires those values.

\subsection{Useful throughput and MFU}

For one causal B1 layer, let $K_m$ be model width and
$D_\Sigma=2D_{QK}+D_V=512$.  The harness counts useful dense and causal
matrix FLOPs as
\begin{equation}
  F_{\mathrm{alg}}=
  4SK_mH(D_\Sigma+D_V)
  +3HS^2(D_{QK}+D_V)
  +\mathbf 1_{\mathrm{wgrad}}\,2SK_mH(D_\Sigma+D_V).
  \label{eq:e2e-algorithmic-flops}
\end{equation}
Elementwise softmax, RoPE, scaling, packing, and reductions are excluded from
this numerator.  Both routes therefore have the same useful-work count, so
the MFU multiplier is exactly the measured speedup.  For an absolute common
denominator, Table~\ref{tab:e2e-mfu} uses 2.5~PFLOP/s dense BF16-equivalent
throughput per GPU, derived from NVIDIA's published two-GPU GB200 superchip
specification and its dense-is-half-sparse note \cite{nvidia2026gb200}.  This
is an algorithmic MFU, not a claim that every operation executes in BF16 or
that a mixed FP4/FP8 kernel should be divided by the FP4 peak.

\begin{table}[htbp]
\centering
\caption{Complete FP4-forward plus FP4+FP8-backward boundary on GPU 0.
Medians use 11 rotated samples after five warmups.  Training includes QKV and
output projection weight gradients.}
\label{tab:e2e-mfu}
\scriptsize
\begin{tabular}{rrr l rrr}
\toprule
S & H & $K_m$ & scope & BF16 / lowp (ms) & speedup & MFU BF16 $\to$ lowp \\
\midrule
4096 & 24 & 3072 & activation & 1.9845 / 1.5696 & 1.264$\times$ & 23.37\% $\to$ 29.55\% \\
4096 & 24 & 3072 & training   & 2.2975 / 1.9114 & 1.202$\times$ & 26.92\% $\to$ 32.36\% \\
8192 & 16 & 2048 & activation & 3.0048 / 2.6881 & 1.118$\times$ & 22.87\% $\to$ 25.56\% \\
8192 & 16 & 2048 & training   & 3.3056 / 3.0152 & 1.096$\times$ & 24.95\% $\to$ 27.35\% \\
8192 & 24 & 4096 & activation & 5.0545 / 4.0739 & 1.241$\times$ & 28.55\% $\to$ 35.42\% \\
8192 & 24 & 4096 & training   & 5.8852 / 4.9502 & 1.189$\times$ & 31.53\% $\to$ 37.48\% \\
4096 & 64 & 8192 & activation & 7.1005 / 4.1286 & 1.720$\times$ & 36.78\% $\to$ 63.25\% \\
4096 & 64 & 8192 & training   & 9.1851 / 6.2505 & 1.469$\times$ & 40.40\% $\to$ 59.37\% \\
8192 & 64 & 8192 & activation & 16.7172 / 10.8004 & 1.548$\times$ & 36.17\% $\to$ 55.99\% \\
8192 & 64 & 8192 & training   & 20.8910 / 15.0497 & 1.388$\times$ & 39.47\% $\to$ 54.79\% \\
\bottomrule
\end{tabular}
\end{table}

The activation rows use the unsignaled stacked handoff; the training rows use
the weight-gradient-capable interleaved handoff.  Both use the FP4 FA4
forward.  A BF16-forward hybrid is faster on several shapes, but is not used
in this table because the measured question is the complete FP4-forward
system.  Across the same five shapes, training-scope QKV weight-gradient
cosine is 0.99007--0.99049, dX is 0.98092--0.98181, dQ/dK are
0.95573--0.95926/0.95523--0.95708, and dV is 0.99008--0.99050.

\subsection{Multi-step convergence proxy}

Operator cosine does not establish optimization behavior.  A separate
teacher/student experiment therefore updates Q, K, V, and output-projection
weights for 500 streamed synthetic activation batches.  Both students begin
from the same 20\%-perturbed teacher weights and use AdamW with learning rate
$2\times10^{-5}$, loss scaling $2^{20}$, and identical data order.  The
low-precision student refreshes packed weight operands after every update and
uses the complete FP4 forward and FP4+FP8 backward.  Validation uses a held-out
activation matrix.

The FP4 forward creates a nonzero teacher quantization floor.  Let
$L_{\mathrm{floor}}$ be the native low-precision validation loss obtained by
the exact teacher weights and define reducible excess
\begin{equation}
  E_{\mathrm{lowp}}=L_{\mathrm{lowp}}-L_{\mathrm{floor}}.
  \label{eq:lowp-excess-loss}
\end{equation}
Table~\ref{tab:convergence-proxy} reports both raw loss and progress above that
measured floor.  ``Improvement retained'' is the low-precision absolute loss
decrease divided by the BF16 decrease.  The final column is
$(E_T/E_0)/(L^{16}_T/L^{16}_0)$; values below one mean the reducible
low-precision error decayed by a larger fraction.

\begin{table}[htbp]
\centering
\caption{Five-hundred-step S4096/H24/K3072 teacher/student convergence proxy.}
\label{tab:convergence-proxy}
\scriptsize
\begin{tabular}{r rrrrr}
\toprule
Seed & BF16 $L_0\to L_T$ & lowp $L_0\to L_T$ & floor & improvement retained & normalized remaining ratio \\
\midrule
2026081417 & 0.08143 $\to$ 0.05766 & 0.14579 $\to$ 0.12212 & 0.09760 & 0.996 & 0.718 \\
2026081431 & 0.08319 $\to$ 0.05750 & 0.14465 $\to$ 0.12141 & 0.09506 & 0.905 & 0.769 \\
2026081447 & 0.08139 $\to$ 0.05675 & 0.14071 $\to$ 0.11861 & 0.09346 & 0.897 & 0.763 \\
\bottomrule
\end{tabular}
\end{table}

The low-precision route retains 93.2\% of BF16's absolute validation-loss
improvement on average (89.7--99.6\%).  Its normalized excess remaining is
0.524 versus 0.699 for BF16's normalized raw loss.  This is evidence that the
mixed backward does not prevent optimization in this block-level test.  It is
not evidence of equal final loss: native low-precision loss remains higher
because of the 0.093--0.098 forward floor, and evaluating the learned
low-precision weights through the common BF16 forward leaves final loss
25.5\% higher on average.  The experiment is not language-model pretraining,
does not include RMSNorm, MLP, residual, tokenizer, or optimizer sharding, and
cannot support a perplexity or stock-Llama convergence claim.

\subsection{Aggressive pure-FP4 projection epilogue}

Table~\ref{tab:pure-single-quant-projection} isolates the final pure-FP4
epilogue change.  ``Dual'' is the earlier unified projection that emitted the
ordinary Q/K layouts and then repeated fixed-scale quantization for the two
pure dQ/dK operands.  ``Single'' uses
Equations~\ref{eq:pure-projection-single-quant}--\ref{eq:pure-projection-layout-fanout}
to publish all six layouts from one code fragment.

\begin{table}[htbp]
\centering
\caption{Unified QKV projection before and after single-quantization pure-Q/K
fanout.  Percentages are reductions in median time.}
\label{tab:pure-single-quant-projection}
\scriptsize
\begin{tabular}{l rrr rrr}
\toprule
Shape & dual + BF16 & single + BF16 & reduction & dual no BF16 & single no BF16 & reduction \\
 & (ms) & (ms) & & (ms) & (ms) & \\
\midrule
S2048/H8/K1024  & 0.107936 & 0.099072 & 8.21\%  & 0.104000 & 0.096896 & 6.83\% \\
S4096/H24/K1024 & 0.446112 & 0.380832 & 14.63\% & 0.410464 & 0.352384 & 14.15\% \\
\bottomrule
\end{tabular}
\end{table}

The composed measurement in Table~\ref{tab:pure-single-quant-chain} charges
projection, Q/K publication, and the pure-MX backward.  The baseline performs
the fixed-scale pure pack in a standalone kernel after projection; the fused
route consumes the projection-native bundle directly.  Prepared-backward
timings remain within about 2~$\mu$s, showing that the gain is removal of
producer-side conversion and movement rather than a changed attention
kernel.

\begin{table}[htbp]
\centering
\caption{Projection plus pure-MX backward with standalone versus
projection-epilogue Q/K publication.}
\label{tab:pure-single-quant-chain}
\small
\begin{tabular}{l rrrr}
\toprule
Shape & separate pack (ms) & fused epilogue (ms) & saving ($\mu$s) & speedup \\
\midrule
S8192/H8/K1024  & 1.102848 & 1.033216 & 69.632 & 1.067$\times$ \\
S4096/H24/K4096 & 1.124192 & 1.027040 & 97.152 & 1.095$\times$ \\
\bottomrule
\end{tabular}
\end{table}

The no-BF16 pure projection specialization compiles at 128 registers and two
barriers with zero spills; the optional-BF16 specialization uses 168
registers, two barriers, and zero spills.  The downstream attention kernel is
unchanged at 128 registers, sixteen barriers, 231,536 bytes of shared memory,
and zero spills.

\subsection{Upstream dO publication}

The fused upstream producer publishes row-major and transposed-column MXFP4
dO plus the FP32 softmax delta in one pass.  At S8192/H8, five warmups and 31
rotated samples give:

\begin{table}[htbp]
\centering
\caption{Upstream dO setup: three split producers versus one fused producer.}
\label{tab:upstream-dout}
\small
\begin{tabular}{l rr r}
\toprule
Route & median (ms) & saving (ms) & speedup \\
\midrule
row quantizer + column quantizer + delta & 0.132352 & --- & --- \\
fused row + column + delta               & 0.068416 & 0.063936 & 1.935$\times$ \\
\bottomrule
\end{tabular}
\end{table}

The four payload/scale outputs are bit-exact between routes.  The fused delta
has maximum absolute difference $5.96\times10^{-8}$ and relative L2 error
$6.72\times10^{-8}$ against the FP32 reference.  This is a producer-side win;
it does not by itself make the standalone all-MX backward the selected route.

\subsection{Represented probability reuse}

Reusing the already-produced represented E4M3 P inside the retained backward
reduces S8192/H8 from 0.892480 to 0.885696 ms, a 6.784 $\mu$s or 0.76\%
saving.  The outputs agree with the recompute route at cosine 0.999004,
0.998987, and 0.999844 for dQ, dK, and dV.  Persisting a complete forward P
cache is rejected because its S8192/H8 payload and scale pages total about
272 MiB.

A stronger three-command dV candidate reused MXFP4 P directly for dS.  It
measured 0.782464 versus 0.756192 ms for the retained winner and returned a
non-finite dV in the comparison.  After reverting its two compile-time gates,
the stable mixed route again produced finite outputs.  The retained endpoint
therefore keeps local E4M3 P reuse and does not advertise cross-forward P
caching or MXFP4 P-to-dS replay.

\subsection{Projection-backward chain}

Writing fully reduced BF16 gradients into one interleaved
$[dQ,dK,dV]$ buffer and selecting a head-contiguous GEMM split gives a smaller
but consistent systems gain:

\begin{table}[htbp]
\centering
\caption{Old three-output attention plus three projection GEMMs versus the
integrated BF16 QKV buffer and auto-split projection dgrad.}
\label{tab:projection-backward-chain}
\small
\begin{tabular}{l rrr}
\toprule
Shape & old chain (ms) & integrated chain (ms) & speedup \\
\midrule
S4096/H64  & 3.276608  & 3.142528  & 1.043$\times$ \\
S8192/H8   & 0.844672  & 0.827136  & 1.021$\times$ \\
S8192/H24  & 2.407808  & 2.391808  & 1.007$\times$ \\
S16384/H64 & 24.055679 & 23.471905 & 1.025$\times$ \\
\bottomrule
\end{tabular}
\end{table}

dK and dV are bit-identical to the direct-BF16 endpoint.  dQ differs by at
most $9.54\times10^{-7}$ from reduction ordering, and the resulting dX cosine
is approximately 0.999992--0.999996.  This path improves materialization and
GEMM geometry, but it does not capture the no-publication floor because global
BF16 dQ reduction remains mandatory.

The head-pipelined consumer supersedes that original whole-row polling
experiment.  Table~\ref{tab:head-pipelined-projection} uses seven warmups
and 31 rotated samples.  The H24 forced row is the diagnostic topology; the
public H24 row is the shape-dispatched fallback.  Positive delta means that
the public or forced entry point is faster than attention plus cuBLAS.

\begin{table}[htbp]
\centering
\caption{Projection backward after moving per-head readiness into the
projection K loop and adding a non-regressing public dispatch.}
\label{tab:head-pipelined-projection}
\small
\begin{tabular}{l l rrrr}
\toprule
Shape & route & attention only & attention + cuBLAS & result & delta \\
 & & (ms) & (ms) & (ms) & ($\mu$s) \\
\midrule
S8192/H8  & public pipelined & 0.746528 & 0.774944 & 0.772352 & +2.592 \\
S4096/H24 & public fallback  & 0.628768 & 0.734336 & 0.733728 & +0.608 \\
S4096/H24 & forced pipelined & 0.633952 & 0.740384 & 0.748256 & -7.872 \\
\bottomrule
\end{tabular}
\end{table}

At H8, moving the wait to the head boundary, adding the async-proxy
publication fence, and alternating two TMEM accumulators turns the old
14--17~$\mu$s loss into a 2.6~$\mu$s win.  One representative invocation had
dX cosine 0.99999535 and relative L2 0.003051; dK and dV remain exact.  At
H24, a four-row-block persistent prefix plus a disjoint cuBLAS suffix reduces
the old 93.6~$\mu$s loss to 7.9~$\mu$s, but does not win.  H16 and H64
similarly favor the ordinary vendor-GEMM chain.  The public entry point
consequently retains the pipelined topology only for H8 and shape-dispatches
wider heads to the non-regressing BF16-dQ-plus-cuBLAS path.  On the final H24
check that fallback was bit-exact to the comparison output and 0.608~$\mu$s
faster, which is timing noise rather than a claimed kernel gain.  The forced
experimental route is available with
\code{TK\_FA4\_DQ\_PROJECTION\_FORCE\_PIPELINED}.

An isolated fixed-dQ test makes the remaining accuracy boundary explicit:
the persistent projection consumer is bit-exact against \code{torch.mm}.
Run-to-run variation originates in the nondeterministic BF16 TMA reduction
order of the transposed attention grid, not in the projection MMA.  Most
single-run cosine values are above 0.9999, but the S8192/H8 repeat screen saw
one 0.99863 cosine / 0.0523 relative-L2 excursion; this is another reason not
to extend the topology to wider heads without changing the reduction owner.

\subsection{Compact dQ projection format screen}

Table~\ref{tab:dq-projection-format} compares tensor formats for the isolated
$dQW_Q$ projection.  The GEMM columns use already-prepared operands; the
charged columns include activation conversion.  Weights remain persistent.
``NV delayed'' uses the device-resident $\bar g_D$ in
Equation~\ref{eq:compact-dq-nvfp4} and therefore skips a matrix-wide amax.

\begin{table}[htbp]
\centering
\caption{dQ projection format timing.  All entries are median milliseconds.}
\label{tab:dq-projection-format}
\scriptsize
\begin{tabular}{l rrrr rrr}
\toprule
Shape & BF16 & FP8 & NVFP4 & MXFP4 & FP8 charged & NV delayed & MX charged \\
\midrule
S8192/H8  & 0.078816 & 0.107776 & 0.045920 & 0.098016 & 0.251392 & 0.061952 & 0.269568 \\
S4096/H24 & 0.155968 & 0.144288 & 0.072192 & 0.125184 & 0.384608 & 0.092192 & 0.306496 \\
\bottomrule
\end{tabular}
\end{table}

The same-step full NVFP4 conversion plus GEMM is 0.070208 and 0.101920 ms on
the two rows.  Projection cosine against BF16 is 0.999301/0.999296 for FP8,
0.991005/0.990976 for NVFP4, and 0.983365/0.983188 for MXFP4.  NVFP4 is the
only FP4 family in this screen that repays its activation conversion while
retaining approximately 0.991 cosine.

The composed chain in Table~\ref{tab:compact-dq-chain} includes attention,
private BF16 dQ reduction, delayed-scale NVFP4 packing, and the projection.
dK and dV are bit-exact with the BF16-chain outputs.

\begin{table}[htbp]
\centering
\caption{Adaptive backward plus dQ projection.  Speedup is BF16 chain divided
by the compact NVFP4 chain.}
\label{tab:compact-dq-chain}
\small
\begin{tabular}{l rrrr}
\toprule
Shape & attention only & BF16 chain & compact NVFP4 chain & speedup / dX cosine \\
\midrule
S8192/H8  & 0.750848 & 0.777568 & 0.792736 & 0.981$\times$ / 0.991005 \\
S4096/H24 & 0.630816 & 0.731968 & 0.705760 & 1.037$\times$ / 0.990976 \\
\bottomrule
\end{tabular}
\end{table}

The H24 projection is wide enough for NVFP4 throughput to dominate and gains
3.7\%.  H8 loses 1.9\% because the smaller projection does not amortize the
extra launch, synchronization, and packing boundary.  A systems dispatcher
should therefore use NVFP4 for the wide dQ projection and retain the existing
BF16 projection path for the shallow H8 case.  Quantizing a completed private
BF16 tensor is a measured intermediate rung; only tile-ready producer fusion
can remove the remaining round trip.

\subsection{Lifetime-safe adaptive mixed dP}

Separating the fixed mixed-dP scale word from the adaptive score scale page
removes the earlier lifetime violation without increasing the 128-register
footprint.  A two-seed, three-regime screen gives:

\begin{table}[htbp]
\centering
\caption{Adaptive FP8 dP/dV versus lifetime-safe adaptive mixed dP.  Timing
ranges span two seeds; component-cosine ranges combine dQ and dK.}
\label{tab:adaptive-mixed}
\scriptsize
\begin{adjustbox}{max width=\textwidth}
\begin{tabular}{l rr r rrr}
\toprule
Shape & adaptive FP8 (ms) & adaptive mixed (ms) & mixed delta & calibrated & unit Q/K & sparse outlier \\
\midrule
S4096/H24 & 0.6589--0.6603 & 0.6420--0.6452 & 2.3--2.7\% faster & 0.9833--0.9838 & 0.9605--0.9620 & 0.9540--0.9552 \\
S8192/H8  & 0.7670--0.7709 & 0.7770--0.7813 & 1.3--1.7\% slower & 0.9827--0.9842 & 0.9588--0.9656 & 0.9547--0.9554 \\
\bottomrule
\end{tabular}
\end{adjustbox}
\end{table}

The adaptive FP8 route remains near 0.9914 calibrated dQ/dK cosine, about
0.978 in the unit-scale regime, and about 0.963 for sparse outliers.  Adaptive
mixed therefore restores correctness and gives a real H24 speed point, but it
does not preserve the retained route's component accuracy and it regresses at
H8.  The public API exposes it explicitly as \code{route="mixed"}; it is not
the automatic default.

\subsection{Faster fixed mixed route}

The fixed-scale mixed E4M3/MXFP4 dP route remains a useful Pareto point.  On
the eight-shape broad matrix it is 13.3--15.6\% faster than BF16 and
3.3--5.3\% faster than the higher-accuracy FP4+FP8 winner.  Its calibrated
dQ/dK cosine is only 0.9832--0.9840, versus 0.9912--0.9915 for the winner.
The lifetime-safe adaptive specialization reaches essentially the same
component-accuracy class.  It is useful at H24, but the fixed route remains
faster and the retained all-E4M3 dP/dV route remains more accurate.

\subsection{NVFP4/MXFP4 pure-format screen}

The pure route was also rebuilt with K16/E4M3 NVFP4 scale pages for dS, the
compact Q/K operands used by dQ/dK, or both.  Table~\ref{tab:pure-format-screen}
reports the two representative shapes.  ``MX/MX'' is the restored pure
baseline; each other row is paired with the BF16 timing from its own binary.

\begin{table}[htbp]
\centering
\caption{Pure-FP4 scale-family combinations.  Format order is dS / compact
Q/K.  Times are milliseconds.}
\label{tab:pure-format-screen}
\scriptsize
\begin{tabular}{l l rr rrr}
\toprule
Shape & formats & BF16 & candidate & dQ cosine & dK cosine & dV cosine \\
\midrule
S8192/H8  & MX/MX & 0.867200 & 0.779232 & 0.881413 & 0.882029 & 0.985053 \\
            & MX/NV & 0.875680 & 0.802656 & 0.858788 & 0.857532 & 0.985056 \\
            & NV/MX & 0.868640 & 0.904320 & 0.969496 & 0.969156 & 0.985056 \\
            & NV/NV & 0.870624 & 0.979264 & non-finite & non-finite & 0.985056 \\
S4096/H24 & MX/MX & 0.707072 & 0.635040 & 0.864334 & 0.864188 & 0.984754 \\
            & MX/NV & 0.706368 & 0.645248 & 0.844903 & 0.846104 & 0.984521 \\
            & NV/MX & 0.707712 & 0.735040 & 0.969362 & 0.969550 & 0.984521 \\
            & NV/NV & 0.708480 & 0.783168 & non-finite & non-finite & 0.984521 \\
\bottomrule
\end{tabular}
\end{table}

NVFP4 dS is an informative accuracy rung: keeping MXFP4 Q/K raises dQ/dK
cosine to approximately 0.969.  Its K16 scale production and movement make it
3.9--4.1\% slower than BF16, however.  NVFP4 Q/K is slower and less accurate
than MXFP4 Q/K in the current layout; enabling both NVFP4 changes is non-finite
and 10.5--12.5\% slower than BF16.  All three experimental gates are restored
to zero.  No NVFP4/MXFP4 combination displaces the hybrid frontier.

\subsection{No complete CuTe native-FP4 comparison}

The local CuTe native-FP4 dK path does not complete in this checkout.  The
runnable dQ-only and BF16 controls have different precision and completeness
contracts.  Accordingly, this report does not quote a native-FP4
ThunderKittens/CuTe speed ratio; doing so would compare a complete kernel with
an incomplete lower bound.
